%\documentclass[hyperref={pdfpagelabels=false},slidetop,9pt]{beamer}
\documentclass[slidetop,8pt]{beamer}
\usepackage[T1]{fontenc}
\usepackage[utf8]{inputenc}
\newcommand{\nom}{Porte conteneur}
\newcommand{\sequence}{03}
\newcommand{\num}{04}
\newcommand{\type}{TD}
\newcommand{\descrip}{Résolution d'un problème en utilisant des méthodes algorithmiques}
\newcommand{\competences}{Alt-C3: Concevoir un algorithme répondant à un problème précisément posé}
\usepackage{etex}
\usepackage{tikz}
\usepackage[european]{circuitikz}
\usepackage{pgf}
\usepackage[all]{xy}
\usepackage{pgfpages}
\usepackage{graphbox}
\usepackage{pdfpages}
\usepackage[adobe-utopia]{mathdesign}
\usepackage{ifthen}
\usepackage{cancel}
\usepackage{framed}
\usepackage{subfig}
\usepackage{tabularx}
\usepackage{setspace}
\usepackage{soul}
\usepackage{schemabloc}
\usepackage{eqnarray}
\usepackage[dot, phantomtext]{dashundergaps}
\usepackage{media9}
\usepackage{multimedia}

\author{Renaud Costadoat}
\institute{Lycée Dorian}

\usepackage{multido}
\usepackage{multirow}
\usepackage{multicol} % Portions de texte en colonnes
\usepackage{flafter}%floatants après la référence

\usepackage{color}
\usepackage{xcolor}
\usepackage{colortbl}

\usepackage[gen]{eurosym}
\usepackage{tikz}
%\usepackage{pstricks,pst-node,pst-tree,pst-solides3d}
\usepackage{lmodern}
\usepackage[francais]{babel}
\usepackage{pslatex}
\usetheme{renaud}
\usepackage{times}
\usepackage{amsmath}
\usepackage{verbatim}
\usepackage{moreverb}
%\usetikzlibrary{arrows,shapes}
\usepackage{graphicx}
\usepackage{psfrag}
\usepackage{wrapfig}
\usepackage{etoolbox}

\definecolor{gris25}{gray}{0.75}
\definecolor{bleu}{RGB}{18,33,98}
\definecolor{bleuf}{RGB}{42,94,171}
\definecolor{bleuc}{RGB}{231,239,247}
\definecolor{rougef}{RGB}{185,18,27}
\definecolor{rougec}{RGB}{255,188,204}%255,230,231
\definecolor{vertf}{RGB}{103,126,82}
\definecolor{vertc}{RGB}{220,255,191}

\setlength\parindent{24pt}
\parskip 7.2pt
\parindent 8pt

\newenvironment{rem}[1][\hsize]%
{%
    \def\FrameCommand
   {%
\rotatebox{90}{\textit{\textsf{Remarque}}} 
       {\color{bleuf}\vrule width 3pt}%
       \hspace{0pt}%must no space.
       \fboxsep=\FrameSep\colorbox{bleuc}%
  }%
    \MakeFramed{\hsize#1\advance\hsize-\width\FrameRestore}%
}%
{\endMakeFramed}%


\newenvironment{savoir}[1][\hsize]%
{%
    \def\FrameCommand
    {%
\rotatebox{90}{\textit{\textsf{Savoir}}} 
        {\color{bleuf}\vrule width 3pt}%
        \hspace{0pt}%must no space.
        \fboxsep=\FrameSep\colorbox{bleuc}%
    }%
    \MakeFramed{\hsize#1\advance\hsize-\width\FrameRestore}%
}%
{\endMakeFramed}%

\newenvironment{prob}[1][\hsize]%
{%
    \def\FrameCommand%
    {%
\rotatebox{90}{\textit{\textsf{Problématique}}} 
        {\color{rougef}\vrule width 3pt}%
        \hspace{0pt}%must no space.
        \fboxsep=\FrameSep\colorbox{rougec}%
    }%
    \MakeFramed{\hsize#1\advance\hsize-\width\FrameRestore}%
}%
{\endMakeFramed}%

\newenvironment{obj}[1][\hsize]%
{%
    \def\FrameCommand%
    {%
\rotatebox{90}{\textit{\textsf{Objectif}}} 
        {\color{vertf}\vrule width 3pt}%
        \hspace{0pt}%must no space.
        \fboxsep=\FrameSep\colorbox{vertc}%
    }%
    \MakeFramed{\hsize#1\advance\hsize-\width\FrameRestore}%
}%
{\endMakeFramed}%

\newenvironment{defi}[1][\hsize]%
{%
    \def\FrameCommand%
    {%
\rotatebox{90}{\textit{\textsf{Definition}}} 
        {\color{bleuf}\vrule width 3pt}%
        \hspace{0pt}%must no space.
        \fboxsep=\FrameSep\colorbox{rougec}%
    }%
    \MakeFramed{\hsize#1\advance\hsize-\width\FrameRestore}%
}%
{\endMakeFramed}%


\newenvironment{hypo}[1][\hsize]%
{%
    \def\FrameCommand%
    {%
\rotatebox{90}{\textit{\textsf{Hypothèse\\}}} 
        {\color{bleuf}\vrule width 3pt}%
        \hspace{0pt}%must no space.
        \fboxsep=\FrameSep\colorbox{bleuc}%
    }%
    \MakeFramed{\hsize#1\advance\hsize-\width\FrameRestore}%
}%
{\endMakeFramed}%


\newenvironment{prop}[1][\hsize]%
{%
    \def\FrameCommand%
    {%
\rotatebox{90}{\textit{\textsf{Propriété}}} 
        {\color{bleuf}\vrule width 3pt}%
        \hspace{0pt}%must no space.
        \fboxsep=\FrameSep\colorbox{bleuc}%
    }%
    \MakeFramed{\hsize#1\advance\hsize-\width\FrameRestore}%
}%
{\endMakeFramed}%

\newenvironment{props}[1][\hsize]%
{%
    \def\FrameCommand%
    {%
\rotatebox{90}{\textit{\textsf{Propriétés}}} 
        {\color{bleuf}\vrule width 3pt}%
        \hspace{0pt}%must no space.
        \fboxsep=\FrameSep\colorbox{bleuc}%
    }%
    \MakeFramed{\hsize#1\advance\hsize-\width\FrameRestore}%
}%
{\endMakeFramed}%

\newenvironment{exemple}[1][\hsize]%
{%
    \def\FrameCommand%
    {%
\rotatebox{90}{\textit{\textsf{Exemple}}} 
        {\color{vertf}\vrule width 3pt}%
        \hspace{0pt}%must no space.
        \fboxsep=\FrameSep\colorbox{vertc}%
    }%
    \MakeFramed{\hsize#1\advance\hsize-\width\FrameRestore}%
}%
{\endMakeFramed}%

\newenvironment{resultat}[1][\hsize]%
{%
    \def\FrameCommand%
    {%
\rotatebox{90}{\textit{\textsf{Resultat}}} 
        {\color{rougef}\vrule width 3pt}%
%        {\color{bleuf}\vrule width 3pt}%
        \hspace{0pt}%must no space.
        \fboxsep=\FrameSep\colorbox{rougec}%
    }%
    \MakeFramed{\hsize#1\advance\hsize-\width\FrameRestore}%
}%
{\endMakeFramed}%

\newenvironment{methode}[1][\hsize]%
{%
    \def\FrameCommand%
    {%
\rotatebox{90}{\textit{\textsf{Méthode\\}}} 
        {\color{rougef}\vrule width 3pt}%
        \hspace{0pt}%must no space.
        \fboxsep=\FrameSep\colorbox{rougec}%
    }%
    \MakeFramed{\hsize#1\advance\hsize-\width\FrameRestore}%
}%
{\endMakeFramed}%

\newenvironment{theo}[1][\hsize]%
{%
    \def\FrameCommand%
    {%
\rotatebox{90}{\textit{\textsf{Théorème\\}}} 
        {\color{rougef}\vrule width 3pt}%
        \hspace{0pt}%must no space.
        \fboxsep=\FrameSep\colorbox{rougec}%
    }%
    \MakeFramed{\hsize#1\advance\hsize-\width\FrameRestore}%
}%
{\endMakeFramed}%

\newenvironment{warn}[1][\hsize]%
{%
    \def\FrameCommand%
    {%
\rotatebox{90}{\textit{\textsf{Attention\\}}} 
        {\color{rougef}\vrule width 3pt}%
        \hspace{0pt}%must no space.
        \fboxsep=\FrameSep\colorbox{rougec}%
    }%
    \MakeFramed{\hsize#1\advance\hsize-\width\FrameRestore}%
}%
{\endMakeFramed}%

% \usepackage{pstricks}
%\usepackage{minitoc}
% \setcounter{minitocdepth}{4}

\setcounter{tocdepth}{2}

% \mtcselectlanguage{french} 

%\usepackage{draftcopy}% "Brouillon"
% \usepackage{floatflt}
\usepackage{psfrag}
%\usepackage{listings} % Permet d'insérer du code de programmation
\renewcommand{\baselinestretch}{1.2}

% Changer la numérotation des figures :
% ------------------------------------
% \makeatletter
% \renewcommand{\thefigure}{\ifnum \c@section>\z@ \thesection.\fi
%  \@arabic\c@figure}
% \@addtoreset{figure}{section}
% \makeatother
 


%%%%%%%%%%%%
% Définition des vecteurs %
%%%%%%%%%%%%
 \newcommand{\vect}[1]{\overrightarrow{#1}}

%%%%%%%%%%%%
% Définition des torseusr %
%%%%%%%%%%%%

 \newcommand{\torseur}[1]{%
\left\{{#1}\right\}
}

\newcommand{\torseurcin}[3]{%
\left\{\mathcal{#1} \left(#2/#3 \right) \right\}
}

\newcommand{\torseurstat}[3]{%
\left\{\mathcal{#1} \left(#2\rightarrow #3 \right) \right\}
}

 \newcommand{\torseurc}[8]{%
%\left\{#1 \right\}=
\left\{
{#1}
\right\}
 = 
\left\{%
\begin{array}{cc}%
{#2} & {#5}\\%
{#3} & {#6}\\%
{#4} & {#7}\\%
\end{array}%
\right\}_{#8}%
}

 \newcommand{\torseurcol}[7]{
\left\{%
\begin{array}{cc}%
{#1} & {#4}\\%
{#2} & {#5}\\%
{#3} & {#6}\\%
\end{array}%
\right\}_{#7}%
}

 \newcommand{\torseurl}[3]{%
%\left\{\mathcal{#1}\right\}_{#2}=%
\left\{%
\begin{array}{l}%
{#1} \\%
{#2} %
\end{array}%
\right\}_{#3}%
}

 \newcommand{\vectv}[3]{%
\vect{V\left( {#1} \in {#2}/{#3}\right)}
}


\newcommand{\vectf}[2]{%
\vect{R\left( {#1} \rightarrow {#2}\right)}
}

\newcommand{\vectm}[3]{%
\vect{\mathcal{M}\left( {#1}, {#2} \rightarrow {#3}\right)}
}


 \newcommand{\vectg}[3]{%
\vect{\Gamma \left( {#1} \in {#2}/{#3}\right)}
}

 \newcommand{\vecto}[2]{%
\vect{\Omega\left( {#1}/{#2}\right)}
}

\newcommand{\reponse}[1][4]
{
\multido{}{#1}
{
\begin{center}
\makebox[0.9\linewidth]{\dotfill} \end{center}
}}


% }$$\left\{\mathcal{#1} \right\}_{#2} =%
% \left\{%
% \begin{array}{c}%
%  #3 \\%
%  #4 %
% \end{array}%
% \right\}_{#5}}


%  ------------------------------------------
% | Modification du formatage des sections : | 
%  ------------------------------------------

% Grands titres :
% ---------------

\newcommand{\titre}[1]{%
\begin{center}
      \bigskip
      \rule{\textwidth}{1pt}
      \par\vspace{0.1cm}
      
      \textbf{\large #1}
      \par\rule{\textwidth}{1pt}
    \end{center}
    \bigskip
  }

% Supprime le numéro du chapitre dans la numérotation des sections:
% -----------------------------------------------------------------
\makeatletter
\renewcommand{\thesection}{\@arabic\c@section}
\makeatother


% \titleformat{\chapter}[display]
% {\normalfont\Large\filcenter}
% {}
% {1pc}
% {\titlerule[1pt]
%   \vspace{1pc}%
%   \Huge}[\vspace{1ex}%
% \titlerule]


%%%% Chapitres Comme PY Pechard %%%%%%%%%
% numéro du chapitre
\DeclareFixedFont{\chapnumfont}{OT1}{phv}{b}{n}{80pt}
% pour le mot « Chapitre »
\DeclareFixedFont{\chapchapfont}{OT1}{phv}{m}{it}{40pt}
% pour le titre
\DeclareFixedFont{\chaptitfont}{T1}{phv}{b}{n}{25pt}

\definecolor{gris}{gray}{0.75}
\setbeamertemplate{section in toc}[sections numbered]

\newlength{\RoundedBoxWidth}
\newsavebox{\GrayRoundedBox}
\newenvironment{GrayBox}[1][\dimexpr\textwidth-4.5ex]%
   {\setlength{\RoundedBoxWidth}{\dimexpr#1}
    \begin{lrbox}{\GrayRoundedBox}
       \begin{minipage}{\RoundedBoxWidth}}%
   {   \end{minipage}
    \end{lrbox}
    \begin{center}
    \begin{tikzpicture}%
       \draw node[draw=bleuf,fill=bleuc,rounded corners,%
             inner sep=2ex,text width=\RoundedBoxWidth]%
             {\usebox{\GrayRoundedBox}};
    \end{tikzpicture}
    \end{center}}
    
\ifdef{\prive}{\pgfpagesuselayout{2 on 1}[a4paper,border shrink=0mm]}
\ifdef{\prive}{\setbeamertemplate{navigation symbols}{}}
\setbeamertemplate{itemize item}[ball]
%\setbeamertemplate{blocks}[rounded]%[shadow=true]
\setbeamercolor{block title}{fg=white,bg=grisf}        % titre block normal 
\setbeamercolor{block body}{fg=grisf,bg=grisc!50}      % corps block normal
\setbeamercolor{block body alerted}{fg=white,bg=warning}   % idem pour un block alerte

\title{\nom}
\date{S\sequence \ - \type\num}

\begin{document}
\shorthandoff{:!}
\bibliographystyle{abbrvnat-fr}

\usebackgroundtemplate%
{%
    \centering\includegraphics[width=\paperwidth]{../../img/fond2}%
}

{
\setbeamertemplate{navigation symbols}{}
\setbeamertemplate{headline}[pagetitre]
\setbeamertemplate{footline}[pagetitre]
\usebackgroundtemplate{\centering\includegraphics[width=\paperwidth]{../../img/fond}}
\frame{\titlepage}
}



\section{Cas triangulaire} 

\begin{frame}[fragile]
\frametitle{Système d'équations}

Soit le système d'équations suivant :
\begin{center}
$\left\{\begin{array}{l}
y_0=a_{0,0}.x_0+a_{0,1}.x_1+a_{0,2}.x_2+...+a_{0,n-1}.x_{n-1} \\
y_1=a_{1,0}.x_1+a_{1,1}.x_2+a_{1,2}.x_3+...+a_{1,n-1}.x_{n-1} \\
...\\
y_{n-1}=a_{n-1,0}.x_0+a_{n-1,1}.x_1+a_{n-1,2}.x_2+...+a_{n-1,n-1}.x_{n-1}
\end{array}\right.$
\end{center}

Ce système d'équations peut s'écrire sous la forme suivante: $Y=A.X$, avec:\\
\begin{center}
$X=\left(\begin{array}{c}
x_0 \\ x_1 \\ ...\\ x_{n-1}
\end{array}\right)$,
$Y=\left(\begin{array}{c}
y_0 \\ y_1 \\ ...\\ y_{n-1}
\end{array}\right)$,
$A=\left(\begin{array}{c c c c c}
a_{0,0} & a_{0,1} & a_{0,2} & ... & a_{0,n-1} \\
a_{1,0} & a_{1,1} & a_{1,2} & ... & a_{1,n-1} \\
... & ... & ... & ... & ... \\
a_{n-1,0} & a_{n-1,1} & a_{n-1,2} & ... & a_{n-1,n-1}
\end{array}\right)$
\end{center}

Ces systèmes sont dit de Cramer, c'est à dire qu'ils admettent une unique solution. La matrice A est donc inversible.

\end{frame}

\begin{frame}[fragile]
\frametitle{Cas triangulaire}

Si la matrice A est triangulaire, la résolution du problème est simple.

Ce système d'équations peut s'écrire sous la forme suivante: $Y=A.X$, avec:\\
\begin{center}
$A=\left(\begin{array}{c c c c c}
a_{0,0} & a_{0,1} & a_{0,2} & ... & a_{0,n-1} \\
0 & a_{1,1} & a_{1,2} & ... & a_{1,n-1} \\
0 & 0 & a_{2,2} & ... & a_{2,n-1} \\
... & ... & ... & ... & ... \\
0 & 0 & 0 & ... & a_{n-1,n-1}
\end{array}\right)$
\end{center}

La dernière équation devient: $y_{n-1}=a_{n-1,n-1}.x_{n-1}$, ainsi $x_{n-1}=\dfrac{y_{n-1}}{a_{n-1,n-1}}$. En remontant tout le système l'ensemble des $x_i$ peuvent être déterminés.

Ainsi, $\forall i \in [0,n-2], x_i=\dfrac{y_i-\sum \limits_{j=i+1}^{n-1} a_{i,j}.x_j}{a_{i,i}}$

\end{frame}

\begin{frame}[fragile]
\frametitle{Cas triangulaire}

Soit la matrice triangulaire $A$ est le vecteur $Y$ suivant:
\begin{center}
$A=\left(\begin{array}{c c c}
2 & 5 & 3 \\
0 & 1 & 2 \\
0 & 0 & 4
\end{array}\right)$,$Y=\left(\begin{array}{c}
6 \\
4 \\
9
\end{array}\right)$
\end{center}

Cela correspond au système d'équations suivant :
\begin{center}
$\left\{\begin{array}{l}
6=2.x_0+5.x_1+3.x_2 \\
4=1.x_1+2.x_2 \\
9=4.x_2
\end{array}\right.$
\end{center}

Cela mène à la résolution suivante :
\begin{center}
$\left\{\begin{array}{l}
x_0=\ifdef{\public}{\uncover<2->{0.875\\}}{ \\}
x_1=\ifdef{\public}{\uncover<2->{-0.5\\}}{ \\}
x_2=\ifdef{\public}{\uncover<2->{2.25}}{}
\end{array}\right.$
\end{center}

\end{frame}

\begin{frame}[fragile]
\frametitle{Cas triangulaire: Code Python}

\begin{GrayBox}[0.85\textwidth]
\begin{semiverbatim}\small
\ifdef{\public}{\uncover<2->{
def triangle(A,b):
    n =len(b)
    x = [0 for i in range(n)]
    x[n-1] = b[n-1]/A[n-1][n-1]
    for i in range(n-2,-1,-1):
        s = 0
        for j in range(i+1,n):
            s = s + A[i][j]*x[j]
        x[i] = (b[i] - s)/ A[i][i]
    return x}}{\vspace{5cm}}
   \end{semiverbatim}
\end{GrayBox}
\end{frame}

\begin{frame}[fragile]
\frametitle{Etape 1: Placer le pivot sur la première ligne}

Soit la matrice $A=\left(\begin{array}{c c c c}
2 & 2 & -3 & 4 \\
-2 & -1 & -3 & 2 \\
6 & 4 & 4 & 5 \\
1 & 5 & 3 & 2
\end{array}\right)$,$Y=\left(\begin{array}{c}
2 \\
-5 \\
16 \\
11
\end{array}\right)$

\begin{enumerate}
 \item Recherche de la plus grande valeur $A[0]_max$ de la première colonne de la matrice: ici $A[2][0]=6$,
 \item Inverser cette ligne et la première de la matrice et faire de même pour le vecteur.
\end{enumerate}

$A=\left(\begin{array}{c c c c}
6 & 4 & 4 & 5 \\
-2 & -1 & -3 & 2 \\
2 & 2 & -3 & 4 \\
1 & 5 & 3 & 2
\end{array}\right)$,$Y=\left(\begin{array}{c}
16 \\
-5 \\
2 \\
11
\end{array}\right)$

\end{frame}

\section{Mise sous forme triangulaire}

\begin{frame}[fragile]
\frametitle{Etape 2: Mettre des 0 sur le reste de la première colonne}

Soit la matrice $A=\left(\begin{array}{c c c c}
6 & 4 & 4 & 5 \\
-2 & -1 & -3 & 2 \\
2 & 2 & -3 & 4 \\
1 & 5 & 3 & 2
\end{array}\right)$,$Y=\left(\begin{array}{c}
16 \\
-5 \\
2 \\
11
\end{array}\right)$

Pour chaque ligne sauf la première
\begin{enumerate}
 \item Calcul du coefficient $x=A[i][0]/A[0][0]$,
 \item Calculer les nouveaux coefficients qui permettent d'annuler le premier:  $A[i][k] = A[i][k] - x*A[j][k]$
\end{enumerate}

$A=\left(\begin{array}{c c c c}
6 & 4 & 4 & 5 \\
\ifdef{\public}{\uncover<2->{0 & 0.33 & -1.66 & 3.66\\}}{ \\}
\ifdef{\public}{\uncover<2->{0 & 0.66 & -4.33 & 2.33\\}}{ \\}
\ifdef{\public}{\uncover<2->{0 & 4.33 & 2.33 & 1.16}}{}
\end{array}\right)$,$Y=\left(\begin{array}{c}
16 \\
0.33 \\
-3.33 \\
8.33
\end{array}\right)$
\end{frame}

\begin{frame}[fragile]
\frametitle{Etape 3: Mettre des 0 sur le reste de la première colonne}

La procédure continue avec la matrice de dimension juste inférieure

$A=\left(\begin{array}{c c c}
0.33 & -1.66 & 3.66 \\
0.66 & -4.33 & 2.33 \\
4.33 & 2.33 & 1.16
\end{array}\right)$,$Y=\left(\begin{array}{c}
0.33 \\
-3.33 \\
8.33
\end{array}\right)$

Le résultat de l'étape de mise sous forme triangulaire est alors:
$A=\left(\begin{array}{c c c c}
6 & 4 & 4 & 5 \\
\ifdef{\public}{\uncover<2->{0 & 4.33 & 2.33 & 1.16\\}}{ \\}
\ifdef{\public}{\uncover<2->{0 & 0 & -4.69 & 2.15\\}}{ \\}
\ifdef{\public}{\uncover<2->{0 & 0 & 0 & 2.73}}{}
\end{array}\right)$,$Y=\left(\begin{array}{c}
16 \\
8.33 \\
-4.61 \\
1.5
\end{array}\right)$

\end{frame}


\begin{frame}[fragile]

Il ne reste plus qu'à utiliser la méthode décrite précédemment afin de déterminer le résultat.

Le résultat final pour ces valeurs est alors: $X=\left(\begin{array}{c}
0.642642642643 \\ 1.10810810811 \\ 1.23723723724 \\ 0.552552552553\end{array}\right)$

\end{frame}

\section{Codage sous Python}

\begin{frame}[fragile]
\frametitle{Etude 1: Code Python}

\textbf{Recherche du pivot et permutation des lignes:}

\begin{GrayBox}[0.85\textwidth]
\begin{semiverbatim}\small
\ifdef{\public}{\uncover<2->{
def recherche_pivot(A, j0):
    n = len(A) 
    imax = j0 
    for i in range(j0+1, n):
        if abs(A[i][j0]) > abs(A[imax][j0]):
            imax = i
    return imax}}{\vspace{5cm}}
   \end{semiverbatim}
\end{GrayBox}

\begin{GrayBox}[0.85\textwidth]
\begin{semiverbatim}\small
\ifdef{\public}{\uncover<2->{
def permutation(A,i,j):
    t = A[i]
    A[i] = A[j]
    A[j] = t}}{\vspace{5cm}}
   \end{semiverbatim}
\end{GrayBox}

\end{frame}

\begin{frame}[fragile]
\frametitle{Etude 2: Code Python}

\textbf{Transvection:}

\begin{GrayBox}[0.85\textwidth]
\begin{semiverbatim}\small
\ifdef{\public}{\uncover<2->{
def transvection(A,i,j,x):
    # si vecteur
    if type(A[0]) != list:
        A[i] = A[i] + x*A[j]
    else:
        n = len(A[0])
        for k in range(n):
            A[i][k] = A[i][k] + x*A[j][k]}}{\vspace{5cm}}
   \end{semiverbatim}
\end{GrayBox}

\end{frame}

\begin{frame}[fragile]
\frametitle{Etude 3: Code Python}

\textbf{Résolution du système:}

\begin{GrayBox}[0.85\textwidth]
\begin{semiverbatim}\small
\ifdef{\public}{\uncover<2->{
def resolution_systeme(A,b):
    n = len(A)

    for i in range(n-1):
        i0 = recherche_pivot(A,i)
        permutation(A,i,i0)
        permutation(b,i,i0) 
        
        for k in range(i+1,n):
            x = A[k][i]/A[i][i]
            transvection(A,k,i,-x)
            transvection(b,k,i,-x) 
            
    return triangle(A,b)}}{\vspace{5cm}}
   \end{semiverbatim}
\end{GrayBox}
\end{frame}


\end{document}
